\subsection{等变 Fukaya 范畴}\label{sec:14e}
现在把讨论的两条线索结合起来。令 $M$ 为带角正合辛流形，配备满足\stref{假设}{14.3}的对合 $\iota_M$，并带有不变的二次复体积形式 $\eta_M^2$。

考虑 $M$ 上预等变正合 Lagrangian 膜，并附加如下条件：若其等变，则须满足\stref{假设}{14.4}；除非它们与不动点集不交，否则这还包含一个法向角的选择。为任意一对这样的膜选择 Floer 数据，要求其关于 $\Z/2$ 作用 $\epsilon$ 自然，即
\begin{equation}\label{eq:14.14}
(H_{L_{0,-}^\#,L_{1,-}^\#},J_{L_{0,-}^\#,L_{1,-}^\#})
=(\iota_M^*H_{L_{0,+}^\#,L_{1,+}^\#},\iota_M^*J_{L_{0,+}^\#,L_{1,+}^\#}).
\end{equation}
\SourcePage{217}
若所涉及的两个膜均等变，则意味着 Floer 数据须 $\iota_M$ 不变；在这种情况下，还按\stref{假设}{14.5}限制选择。也要求所有 Floer 数据正则。对等变膜对，\stref{引理}{14.6}之后已说明能够做到；在所有其他情况下也没有困难，因为\eref{14.14}联系的是两个不同的 Floer 数据，只要选择其中一个正则，另一个显然也有同样性质。Floer 数据选择中的对称性意味着存在典范同构
\begin{equation}\label{eq:14.15}
\epsilon:CF^*(L_{0,+}^\#,L_{1,+}^\#)\longrightarrow CF^*(L_{0,-}^\#,L_{1,-}^\#),
\end{equation}
它们与微分严格交换，且平方为恒等。随后，为对象的所有 $(d+1)$ 元组选择扰动数据，满足通常的一致性要求及\eref{14.14}的对应条件。由\stref{引理}{14.7}，可在这一类别内达到正则性，所得复合映射 $\mu^d$ 与\eref{14.15}严格交换。由此得到 $\Ai$ 范畴 $\F(M)^{\mathrm{naive}}$，称为\emph{等变 Fukaya 范畴}；其对象是满足上述条件的预等变膜，态射空间为 $CF^*(L_{0,+}^\#,L_{1,+}^\#)$。此 $\Ai$ 范畴带有由\eref{14.15}给出的朴素 $\Z/2$ 作用。为与之前的构造比较，转到上同调层次；显然存在与严格 $\Z/2$ 作用相容的全忠实嵌入 $H\F(M)^{\mathrm{naive}}\to H\F(M)^{\mathrm{strict}}$。事实上，这是一个等价：不满足\stref{假设}{14.4}的不变 Lagrangian 子流形，不能实现为 $H\F(M)^{\mathrm{naive}}$ 中的等变膜，但可作为预等变膜纳入，而在上同调层次这些对象同构。换言之，这样的子流形作为 $H\F(M)^{\mathrm{naive}}$ 的对象存在，但无法成为 $\Z/2$ 作用的不动点。这样做的好处是，$H\F(M)^{\mathrm{naive}}$ 确实是某个带朴素 $\Z/2$ 作用的 $\Ai$ 范畴的上同调范畴；$H\F(M)^{\mathrm{strict}}$ 则只是上同调层次的构造。\znote{209}{两者的差别只在链层次可见：$\F(M)^{\mathrm{naive}}$ 上的 $\Z/2$ 作用由链映射\eref{14.15}给出，是严格作用；而 $H\F(M)^{\mathrm{strict}}$ 只记住上同调上的作用，无法再提升到链层次。二者的上同调范畴等价，故就“分次范畴”而言无差别；但凡需要在上链层次作等变构造（取商、扭曲复形、等变 Floer 理论等），都必须使用朴素版本，否则会丢失高阶同伦信息。这正是不满足\stref{假设}{14.4}的不变 Lagrangian 子流形在 $H\F(M)^{\mathrm{naive}}$ 中不成为不动点的原因：它们只是同构于不动点，而同构并非链层次的相等。[来源：原书 §(14e)，印刷页 209。]}

为更详细地描述\eref{14.15}，取 $y_+\in\mathcal C(L_{0,+}^\#,L_{1,+}^\#)$，并考虑相关的一对线性 Lagrangian 膜 $\Lambda_{k,+}^\#$，即 $L_{k,+}^\#$ 在 $y_+(k)$ 处的切空间。这里把膜结构加入记号，因为 Hamilton 函数 $H$，因而集合 $\mathcal C$，现在依赖于这些结构。对 $\mathcal C(L_{0,-}^\#,L_{1,-}^\#)$ 中相应的元素 $y_-=\iota_M(y_+)$，由对象上 $\Z/2$ 作用的定义，有关的线性 Lagrangian 膜 $\Lambda_{k,-}^\#$ 是 $L_{k,-}^\#$ 的切空间。作为等变膜结构的一部分，有同构
\begin{equation}\label{eq:14.16}
(f_{k,-})_{y_-(k)}:\Lambda_{k,-}^\#\longrightarrow
\Lambda_{k,+}^\#\otimes\lambda^{\mathrm{top}}(TL_{k,+})_{y_+(k)},
\end{equation}
覆盖 $D\iota_M:(TL_{k,-})_{y_-(k)}\to(TL_{k,+})_{y_+(k)}$ 的作用。由\stref{笔记}{11.18}，有选定的同构 $\lambda^{\mathrm{top}}(TL_{k,+})_{y_+(k)}\cong\xi_{y_+(k)}$，还可沿 $y$ 平行移动，认同 $\xi_{y_+(0)}\cong\xi_{y_+(1)}$。将所涉及的两个线性 Lagrangian 膜同时与同一实直线作张量积，不改变绝对指标，也不改变定向空间，后者的“不改变”是在典范同构意义下而言。因此，\eref{14.16}
\SourcePage{218}
诱导同构 $o(y_-)\to o(y_+)$。将这些同构的逆作 $\K$ 规范化，就组成\eref{14.15}。这构成 $\Z/2$ 作用，归结于对\eref{14.2}的条件，以及下述事实：若用非平凡元素 $-e\in\ker(\operatorname{Pin}_n\to O_n)$ 同时作用于两个线性 Lagrangian 膜 $\Lambda_{k,+}^\#$，则在 $o(y)$ 上诱导的自同构是恒等。

一个也许值得特别关注的情形，是两个膜 $L_k^\#$ 都等变，且 $y\in\mathcal C(L_0^\#,L_1^\#)$ 在 $\iota_M$ 的作用下不变。此时，上述构造给出 $o(y)$ 的自同构，其符号由局部几何数据决定。与\stref{引理}{14.9}的证明一样，作简化假设 $H=0$，故 $y$ 为不变交点。回顾\stref{引理}{14.2}：同构 $f_k$ 为不动部分 $L_k^{\mathrm{fix}}$ 的法丛确定选定的平凡化。经 $D\pi$ 投影，这些平凡化给出直线 $e^{i\sigma_k}\R$ 的定向。若该定向与生成元 $e^{i\sigma_k}$ 所给定向一致，置 $w_k=+1$，否则置 $w_k=-1$。于是，

\StatementNumber{11}\begin{lemma}\label{st:14.11}
若 $w_0w_1=+1$，则 $|o(y)|_{\K}$ 属于 $CF^*(L_0^\#,L_1^\#)$ 的 $\epsilon$ 不变部分；否则属于反不变部分。
\end{lemma}
\begin{proof}
按照\stref{笔记}{14.10}的策略，选取从 $\lambda(0)=(TL_0)_y$ 到 $\lambda(1)=(TL_1)_y$ 的路径 $\lambda\in\mathcal P^-\operatorname{Gr}(V)$，完全由 $\Z/2$ 不变 Lagrangian 子空间组成。记 $F_\lambda$ 为 $[0,1]$ 上的关联向量丛。但这次要稍稍改变反不变直和项的同伦类：要求 $\lambda^-(s)=e^{i\sigma(s)}\R$，其中 $\sigma(0)=\sigma_0$、$\sigma(1)=\sigma_1+k\pi$，$k\in\Z$ 满足 $(-1)^k=w_0w_1$。为 $F_\lambda$ 选择膜结构，由实值相位 $\alpha_\lambda^\#$ 和 Pin 结构 $P_\lambda^\#$ 组成，在端点上限制为给定的线性 Lagrangian 膜。为实现\eref{14.15}，还要有同构
\[
f_\lambda:P_\lambda^\#\longrightarrow P_\lambda^\#\otimes\lambda^{\mathrm{top}}(F_\lambda),
\]
覆盖 $D\iota_M$ 在 $F_\lambda$ 上的作用。\stref{引理}{14.2}表明，这样的同构与 $\lambda^-$ 的定向一一对应。显然可调整 $f_\lambda$，使其在 $s=0$ 时与\eref{14.16}一致；追踪诱导的定向，并利用关于 $\lambda^-$ 同伦类的假设，便知在 $s=1$ 时也一致。因此，可在\eref{11.25}中考察 $D\iota_M$ 对\eref{11.20}的作用，读出 $o(y)$ 上诱导的自同构。在此，对每个 $s<1$ 时的非零交 $\lambda^-(s)\cap\lambda^-(1)$，都会产生一个 $(-1)$ 因子。这种点的个数的奇偶性由 $w_0w_1=(-1)^k$ 确定，即与 $k$ 相同，因而得到所述结论。等价地，也可如\stref{引理}{14.9}的证明，将此论证转回指标理论的语言。
\end{proof}
